\documentclass[11pt]{article}
\usepackage[margin=1in]{geometry}
\usepackage[T1]{fontenc}
\usepackage[utf8]{inputenc}
\usepackage{amsmath}

\title{Question 1 Search Answers}
\date{}

\begin{document}

\maketitle

The problem states that the path cost is the number of moves made to reach the current position, so each knight move has unit cost and the optimal cost from $(0,0)$ to $(0,4)$ is $2$.

\begin{enumerate}
\item \textbf{Maximum branching factor:} $8$.

The knight has $8$ legal moves on an infinite chessboard.

\item \textbf{Search techniques guaranteed to find a solution with infinite time:} Breadth-First Search and Iterative-Deepening Depth-First Search.

Depth-First Search is not guaranteed to find a solution on an infinite search tree because it can follow an infinite branch forever.

\item \textbf{Heuristics that allow A* to find an optimal solution of cost $2$:} $h_1$, $h_2$, and $h_3$.

\begin{itemize}
\item $h_1(n)=C$ gives $f(n)=g(n)+C$, which preserves the same ordering as $g(n)$.
\item $h_2(n)=C-g(n)$ gives $f(n)=C$, so A* reduces to FIFO queue behavior, which is Breadth-First Search here.
\item $h_3(n)=2g(n)$ gives $f(n)=3g(n)$, which also preserves the same ordering as $g(n)$.
\item $h_4(n)=-2g(n)$ gives $f(n)=-g(n)$, which prefers deeper nodes and is not guaranteed to return the optimal depth-$2$ solution first.
\end{itemize}

\item \textbf{Heuristics that make A* behave like Breadth-First Search in this problem:} $h_1$, $h_2$, and $h_3$.

Because every move has unit cost, Breadth-First Search expands by depth. The order induced by $f(n)=g(n)+C$, $f(n)=C$, and $f(n)=3g(n)$ matches Breadth-First Search under the stated FIFO tie-breaking.

\item \textbf{Heuristics that make A* behave like Uniform-Cost Search in this problem:} $h_1$, $h_2$, and $h_3$.

In this problem Uniform-Cost Search and Breadth-First Search behave the same because every move has cost $1$. Therefore the same three heuristics from Question 4 also match Uniform-Cost Search behavior.

\item \textbf{Consistent heuristics:} $h_1$ and $h_2$.

For consistency we require
\[
h(n) \le c(n,a,n') + h(n').
\]
Since each move has cost $1$:
\begin{itemize}
\item For $h_1(n)=C$, we get $C \le 1 + C$.
\item For $h_2(n)=C-g(n)$ and $g(n')=g(n)+1$, we get
\[
C-g(n) \le 1 + C - g(n') = 1 + C - (g(n)+1) = C-g(n).
\]
So equality holds for every successor.
\end{itemize}

\item \textbf{Landmark heuristic value:} $6$.

Using
\[
h_L(n)=\min_{L \in \text{Landmarks}} C^*(n,L) + C^*(L,\text{goal}),
\]
for $n=(4,-2)$ and goal $(0,4)$:
\begin{align*}
(0,8):\;& 6 + 2 = 8,\\
(8,0):\;& 2 + 4 = 6,\\
(0,-8):\;& 4 + 6 = 10,\\
(-8,0):\;& 6 + 4 = 10.
\end{align*}
Therefore,
\[
h_L(n)=6.
\]
\end{enumerate}

\end{document}